\documentclass[11pt,a4paper]{article}

\usepackage{ctex}
\usepackage[T1]{fontenc}
\usepackage[top=2.5cm, bottom=2.5cm, left=2.5cm, right=2.5cm]{geometry}
\usepackage{booktabs}
\usepackage{array}
\usepackage{enumitem}
\usepackage{tabularx}
\usepackage{amssymb}
\usepackage{xcolor}
\definecolor{primary}{HTML}{1a1a2e}
\definecolor{accent}{HTML}{3b82f6}
\definecolor{success}{HTML}{22c55e}
\definecolor{danger}{HTML}{ef4444}
\definecolor{warning}{HTML}{f59e0b}
\definecolor{graytext}{HTML}{6b7280}
\definecolor{progress}{HTML}{8b5cf6}
\usepackage[hidelinks]{hyperref}
\hypersetup{colorlinks=true, linkcolor=primary, urlcolor=accent}
\usepackage{titlesec}
\titleformat{\section}{\Large\bfseries\color{primary}}{}{0em}{}[\titlerule]
\titleformat{\subsection}{\large\bfseries\color{primary}}{}{0em}{}
\usepackage{fancyhdr}
\pagestyle{fancy}
\fancyhf{}
\fancyhead[L]{\small\color{graytext} 企业级 AI Agent 应用商店 -- 需求文档}
\fancyhead[R]{\small\color{graytext} DEMAND-003}
\fancyfoot[C]{\small\color{graytext} \thepage}
\renewcommand{\headrulewidth}{0.4pt}
\usepackage{tikz}
\usetikzlibrary{shapes.geometric, arrows.meta, positioning, fit}

\title{\textbf{DEMAND-003：任务流转路径定义}\\[0.5em]\small 魏子政 $\to$ OpenClaw $\to$ Cursor 的完整链路（含拍板机制）}
\author{万能产品小助手（PM AI）\\魏子政 审阅}
\date{2026-05-06 / 版本 v0.2}

\begin{document}

\maketitle

\begin{center}
\begin{tabular}{>{\bfseries}p{3cm} p{10cm}}
\toprule
提出日期 & 2026-05-06 \\
提出者 & 魏子政 \\
优先级 & \textcolor{danger}{P0} \\
状态 & \textcolor{progress}{进行中} \\
总需求编号 & D-009 \\
版本变更 & v0.1 $\to$ v0.2：新增拍板机制，明确工具链仅 OpenClaw + Cursor \\
\bottomrule
\end{tabular}
\end{center}

\vspace{1em}

\section{需求描述}

定义魏子政（甲方）$\to$ OpenClaw（PM AI）$\to$ Cursor（开发 AI）的完整任务流转路径。

\textbf{核心原则}：
\begin{enumerate}[nosep]
  \item \textbf{工具链锁定}：仅使用 OpenClaw + Cursor，Claude/Codex 等暂不参与
  \item \textbf{拍板机制}：任何有歧义的环节，OpenClaw 必须先让魏子政拍板决定方向，再组织下一步的开发/验收/commit
  \item \textbf{单线沟通}：魏子政只与 OpenClaw 对话，不直接与 Cursor 交互
\end{enumerate}

\section{流转总览}

\begin{center}
\begin{tikzpicture}[
  node distance=1.2cm,
  box/.style={rectangle, rounded corners, draw=primary, fill=primary!5, minimum width=2.6cm, minimum height=0.9cm, align=center, font=\small},
  decision/.style={diamond, draw=warning, fill=warning!8, minimum width=2cm, align=center, font=\small, aspect=2.2},
  approve/.style={diamond, draw=danger, fill=danger!8, minimum width=2cm, align=center, font=\small, aspect=2.2},
  arrow/.style={-{Stealth[length=2.5mm]}, thick, graytext}
]
  \node[box] (wei) {魏子政\\（甲方）};
  \node[box, right=1.5cm of wei] (oc) {OpenClaw\\（PM AI）};
  \node[decision, below=1cm of oc] (amb) {有歧义？};
  \node[box, left=1.5cm of amb] (decide) {魏子政拍板\\决策方向};
  \node[box, right=1.5cm of oc] (cur) {Cursor\\（开发 AI）};
  \node[decision, below=1cm of cur] (pass) {PM 验收\\通过？};
  \node[box, below=1cm of pass] (commit) {Git Commit\\归档文档};
  \node[box, below=1cm of commit] (report) {汇报魏子政\\交付完成};

  \draw[arrow] (wei) -- node[above, font=\footnotesize]{提出需求} (oc);
  \draw[arrow] (oc) -- node[above, font=\footnotesize]{需求文档+任务} (cur);
  \draw[arrow] (cur) -- (pass);
  \draw[arrow] (pass) -- node[right, font=\footnotesize]{否} (oc);
  \draw[arrow] (pass) -- node[left, font=\footnotesize]{是} (commit);
  \draw[arrow] (commit) -- (report);
  \draw[arrow] (report) -- (wei);
  \draw[arrow] (oc) -- (amb);
  \draw[arrow] (amb) -- node[above, font=\footnotesize]{是} (decide);
  \draw[arrow] (amb) -- node[right, font=\footnotesize, xshift=0.3cm]{否} ++(1.5,0);
  \draw[arrow] (decide) -- (oc);
\end{tikzpicture}
\end{center}

\section{完整流程（7 步）}

\begin{enumerate}
  \item \textbf{魏子政提出需求}：通过飞书发送自然语言需求给 OpenClaw
  \item \textbf{OpenClaw 分析需求}：理解需求，识别是否有歧义或需要决策的环节
  \item \textbf{拍板环节}（关键新增）：
    \begin{itemize}[nosep]
      \item \textbf{无歧义}：OpenClaw 直接生成需求文档，进入第 4 步
      \item \textbf{有歧义}：OpenClaw 整理出需要决策的问题清单，发送给魏子政拍板
      \item 魏子政拍板后，OpenClaw 按拍板结果生成需求文档
      \item \textbf{拍板触发条件}：技术方案有多个选项、需求描述不明确、优先级有冲突、架构设计有分歧
    \end{itemize}
  \item \textbf{OpenClaw 生成需求文档}：DEMAND-xxx tex/pdf，同步更新 MASTER
  \item \textbf{OpenClaw 调度 Cursor}：通过 \texttt{cursor agent --print --trust --workspace <path>} 下发任务
  \item \textbf{Cursor 执行 + 自验收}：开发代码，运行测试，报告结果
  \item \textbf{OpenClaw 独立验收}：
    \begin{itemize}[nosep]
      \item 运行命令、检查文件、调用 API，不依赖 Cursor 自检结果
      \item 验收通过 $\to$ Git commit + 归档文档 + 汇报魏子政
      \item 验收不通过 $\to$ 记录问题，重新下发 Cursor 修复（回到第 5 步）
    \end{itemize}
\end{enumerate}

\section{各节点分工}

\subsection{魏子政（甲方）}

\begin{table}[h]
\centering
\begin{tabularx}{\linewidth}{>{\bfseries}p{2.5cm} X}
\toprule
角色 & 甲方（需求提出者 + 拍板决策者） \\
\midrule
输入 & 业务想法、功能需求、修改意见、优先级指示 \\
输出 & 自然语言需求描述、拍板决策结果 \\
职责 & 提出需求、拍板决策（有歧义时）、审阅需求文档、确认最终验收 \\
不做 & 不直接与 Cursor 对话、不写代码、不调 API \\
\bottomrule
\end{tabularx}
\end{table}

\subsection{OpenClaw（PM AI）}

\begin{table}[h]
\centering
\begin{tabularx}{\linewidth}{>{\bfseries}p{2.5cm} X}
\toprule
角色 & 产品经理（需求 + 拍板协调 + 拆解 + 验收） \\
\midrule
输入 & 魏子政的需求、魏子政的拍板结果、Cursor 的执行结果 \\
输出 & DEMAND-xxx 需求文档、Cursor 任务指令、验收报告 \\
职责 & \begin{itemize}[nosep]
  \item 需求理解与歧义识别
  \item \textbf{有歧义时整理问题清单，请魏子政拍板}（不自行假设）
  \item 按拍板结果生成需求文档（DEMAND-xxx tex/pdf）
  \item 任务拆解 + 调度 Cursor
  \item 独立验收（不依赖 Cursor 自检）
  \item Git commit + 归档文档 + 同步 MASTER
  \item 汇报魏子政
\end{itemize} \\
不做 & 不写业务代码、不做 UI 设计、不直接改数据库 \\
\bottomrule
\end{tabularx}
\end{table}

\subsection{Cursor（开发 AI）}

\begin{table}[h]
\centering
\begin{tabularx}{\linewidth}{>{\bfseries}p{2.5cm} X}
\toprule
角色 & 全栈开发（前端 + 后端 + 环境 + 运维 + 测试） \\
\midrule
输入 & OpenClaw 下发的任务指令（含需求描述、技术约束、目录规范） \\
输出 & 可运行的代码、测试结果、API 响应、构建产物 \\
职责 & \begin{itemize}[nosep]
  \item 遵循 .cursor/rules/ 下的所有规范文件
  \item 前端开发（Vue 3 + Element Plus，含 UI 设计实现）
  \item 后端开发（FastAPI + SQLAlchemy，含 API 实现）
  \item 数据库操作（Alembic 迁移、种子数据）
  \item 环境搭建（Docker Compose、依赖安装）
  \item 自验收（build 通过、测试通过、API 可用）
\end{itemize} \\
不做 & 不生成需求文档、不修改 MASTER、不与魏子政直接对话 \\
\bottomrule
\end{tabularx}
\end{table}

\section{工具链}

\begin{table}[h]
\centering
\begin{tabularx}{\linewidth}{lXX}
\toprule
\textbf{工具} & \textbf{角色} & \textbf{状态} \\
\midrule
OpenClaw & PM AI（需求 + 拍板协调 + 验收） & \textcolor{success}{使用中} \\
Cursor & 开发 AI（前端 + 后端 + 环境 + 测试） & \textcolor{success}{使用中} \\
Claude & 辅助 AI & \textcolor{danger}{暂不参与} \\
Codex & 备选 AI & \textcolor{danger}{暂不参与} \\
\bottomrule
\end{tabularx}
\end{table}

\section{拍板机制详解}

\subsection{什么情况需要拍板}

\begin{itemize}[nosep]
  \item 技术方案有多个可行选项（如：选 A 库还是 B 库）
  \item 需求描述存在歧义或多种理解方式
  \item 功能优先级有冲突（资源有限时先做哪个）
  \item 架构设计有不同方向（如：同步还是异步）
  \item 涉及外部依赖或第三方服务的选择
  \item 需求变更影响已完成的功能
\end{itemize}

\subsection{拍板流程}

\begin{enumerate}[nosep]
  \item OpenClaw 识别到歧义点
  \item 整理为问题清单，列出每个选项的优缺点
  \item 发送给魏子政，请其拍板
  \item 魏子政决策后，OpenClaw 按决策结果继续
  \item \textbf{禁止 OpenClaw 在未拍板的情况下自行假设方向}
\end{enumerate}

\subsection{无需拍板的情况}

\begin{itemize}[nosep]
  \item 技术栈已锁定（MASTER/RFC/TRD 中已决定）
  \item 目录结构已定义（GUIDE-001）
  \item 编码规范已确定（.mdc 规则文件）
  \item 验收标准已在需求文档中明确
\end{itemize}

\section{质量关卡}

\begin{table}[h]
\centering
\begin{tabularx}{\linewidth}{llX}
\toprule
\textbf{关卡} & \textbf{执行者} & \textbf{通过条件} \\
\midrule
拍板确认 & 魏子政 & 歧义点已决策，方向明确 \\
需求确认 & 魏子政 & DEMAND-xxx 文档准确、任务分解合理 \\
Cursor 自验收 & Cursor & build 通过、测试通过、Swagger 可访问 \\
PM 独立验收 & OpenClaw & 运行命令/检查文件/调 API 验证通过 \\
最终验收 & 魏子政 & 确认交付物满足需求 \\
\bottomrule
\end{tabularx}
\end{table}

\section{异常处理}

\begin{table}[h]
\centering
\begin{tabularx}{\linewidth}{lX}
\toprule
\textbf{异常} & \textbf{处理方式} \\
\midrule
Cursor 执行失败 & OpenClaw 分析日志，调整指令或拆分任务，重新下发 \\
需求理解有歧义 & \textbf{触发拍板机制}，请魏子政决策，不自行假设 \\
技术方案不可行 & OpenClaw 整理替代方案，请魏子政拍板选择 \\
依赖安装失败 & OpenClaw 检查环境/版本冲突，Cursor 重试或换依赖 \\
PM 验收不通过 & OpenClaw 记录问题，重新下发 Cursor 修复 \\
魏子政驳回最终验收 & OpenClaw 根据反馈修改需求文档，重新走流程 \\
\bottomrule
\end{tabularx}
\end{table}

\section{验收标准}

\begin{itemize}
  \item[$\square$] 拍板机制已定义并经魏子政确认
  \item[$\square$] 工具链确认仅 OpenClaw + Cursor
  \item[$\square$] 完整流程可走通：需求提出 $\to$ 拍板（如需）$\to$ 文档生成 $\to$ Cursor 调度 $\to$ 验收 $\to$ commit $\to$ 汇报
\end{itemize}

\section*{更新记录}

\begin{tabular}{lll}
\toprule
版本 & 日期 & 变更内容 \\
\midrule
v0.2 & 2026-05-06 & 新增拍板机制；明确工具链仅 OpenClaw+Cursor；流程从 4 步扩展为 7 步 \\
v0.1 & 2026-05-06 & 初始版本 \\
\bottomrule
\end{tabular}

\end{document}
